\documentclass[../../../include/open-logic-section]{subfiles}

\begin{document}

\olfileid{sth}{replacement}{intro}
\olsection{પરિચય}

પ્રતિસ્થાપન એ સ્વયંસિદ્ધિઓનું પ્રરૂપ છે જે $\ZF$
અને~$\Z$ વચ્ચેનો ભેદ ઊભો કરે છે. આપણે
\crefrange{sth:ordinals::chap}{sth:spine::chap} દરમિયાન તેનો
ઉપયોગ કરી લીધો હતો. આ પ્રકરણમાં છેવટે પ્રશ્ન વિચારશું:
પ્રતિસ્થાપન વાજબી છે ખરું?

પ્રશ્નને વધુ ચોક્કસ બનાવવા નોંધવું જોઈએ કે પ્રતિસ્થાપન ખરેખર
ઘણું \emph{શક્તિશાળી} છે. આ પ્રકરણમાં (અને ફરી
\olref[sth][card-arithmetic][fix]{sec}માં) તે કેટલું શક્તિશાળી
છે તેનો ખ્યાલ આવશે. તેનાથી તેના ઔચિત્યની જરૂર પણ સ્પષ્ટ થશે.

આપણે બે પ્રકારનું ઔચિત્ય ચર્ચીશું. અંદાજે, \emph{બાહ્ય}
ઔચિત્યમાં સ્વયંસિદ્ધિથી મળતાં પરિણામો વડે તેને વાજબી
ઠેરવવાનો પ્રયાસ થાય છે; \emph{આંતરિક} ઔચિત્યમાં વિષયના
ગાણિતિક ખ્યાલો પોતે સ્વયંસિદ્ધિને સમર્થન આપે છે એમ બતાવવાનો
પ્રયાસ થાય છે. આનો અર્થ આ પ્રકરણમાં વધુ સ્પષ્ટ થશે, પણ આ તો
મોટા વિષયનો માત્ર આરંભ છે. વધુ માટે ખાસ કરીને
\citeauthor{Maddy1988a} (\citeyear{Maddy1988a} અને
\citeyear{Maddy1988b}) જુઓ.

\end{document}
